\documentclass[11pt,a4paper,oneside]{amsbook}
\usepackage[T1]{fontenc}
\usepackage[utf8]{inputenc}
\usepackage{lmodern}
\usepackage[american]{babel}
\usepackage{amsmath,amssymb,amsthm,mathtools}
\usepackage[all,cmtip]{xy}
\usepackage{microtype}
\usepackage{enumitem}
\usepackage{xspace}
\usepackage{geometry}
\geometry{margin=1.1in}
\usepackage[hidelinks]{hyperref}

% SGA notation used in the translated body fragments.
\DeclareMathOperator{\Hom}{Hom}
\DeclareMathOperator{\Spec}{Spec}
\DeclareMathOperator{\trace}{tr}
\DeclareMathOperator{\gr}{gr}
\newcommand{\id}{\mathrm{id}}
\newcommand{\kres}{\mathit{k}}
\newcommand{\isomto}{\overset{\sim}{\longrightarrow}}
\newcommand{\isomfrom}{\overset{\sim}{\longleftarrow}}
\newcommand{\ZZ}{\mathbf{Z}}
\newcommand{\RR}{\mathbf{R}}
\newcommand{\CC}{\mathbf{C}}
\let\cal\mathcal
\let\goth\mathfrak
\def\to{\mathchoice{\longrightarrow}{\rightarrow}{\rightarrow}{\rightarrow}}
\def\mto{\mathchoice{\longmapsto}{\mapsto}{\mapsto}{\mapsto}}
\def\To{\mathchoice{\Longrightarrow}{\Rightarrow}{\Rightarrow}{\Rightarrow}}
\def\tbigwedge{\mathop{\textstyle\bigwedge}\limits}
\def\cf{{cf.}\kern.3em}
\def\Cf{{Cf.}\kern.3em}
\def\ie{{i.e.}\kern.3em}
\def\resp{resp.\xspace}
\def\ptbl{.\kern.2em }
\newcommand{\No}{no.~}

% References outside this exposé retain their source keys and
% printed numbers. They become ordinary references when a label is added.
\makeatletter
\newcommand{\SourceRef}[2]{\@ifundefined{r@#1}{#2}{\ref{#1}}}
\makeatother

\addto\captionsamerican{\renewcommand{\chaptername}{Expos\'e}}
\renewcommand{\thechapter}{\Roman{chapter}}
\renewcommand{\thesection}{\arabic{section}}
\theoremstyle{plain}
\newtheorem{theorem}{Theorem}[section]
\newtheorem{proposition}[theorem]{Proposition}
\newtheorem{lemma}[theorem]{Lemma}
\newtheorem{corollary}[theorem]{Corollary}
\newtheorem{corollaries}[theorem]{Corollaries}
\theoremstyle{definition}
\newtheorem{definition}[theorem]{Definition}
\theoremstyle{remark}
\newtheorem*{remarks}{Remarks}
\newtheorem*{remark}{Remark}
\newtheorem*{scholium}{Scholium}
\newenvironment{namedstatement}[1]
  {\par\addvspace{.5\baselineskip}\noindent\textsc{#1.}\ \itshape\ignorespaces}
  {\par\addvspace{.5\baselineskip}}

\title[SGA~1, Expos\'e I (English draft)]{\'Etale morphisms\\
{\normalsize SGA~1, Expos\'e I\\English draft}}
\author{A.\ Grothendieck}
\date{}
\hypersetup{
  pdftitle={SGA 1, Expose I: Etale morphisms (English draft)},
  pdfauthor={A. Grothendieck},
  pdfsubject={Unofficial English translation of SGA 1, Expose I}
}

\begin{document}
\frontmatter
\maketitle

\chapter*{About this draft}
This is an unofficial English translation of Expos\'e~I,
``Morphismes \'etales'', from SGA~1. The source volume is
A.\ Grothendieck and M.\ Raynaud,
\textit{Rev\^etements \'etales et groupe fondamental} (SGA~1),
S\'eminaire de g\'eom\'etrie alg\'ebrique du Bois Marie, 1960--61.
It follows the slightly corrected SMF recomposition,
\href{https://arxiv.org/abs/math/0206203v2}{arXiv:math/0206203v2}.

This draft contains the opening convention and all eleven sections,
including their proofs and footnotes. Statement numbering and
mathematical notation follow the source, including its two statements
numbered~9.2 (a corollary and a proposition).
References to material outside this draft retain their original
numbers. Apparent mathematical misprints in the source have been
retained; they are recorded in the accompanying README, separately
from the translated text. Scholarly proofreading remains outstanding.

The translator's contribution is licensed under
\href{https://creativecommons.org/licenses/by-sa/4.0/}{CC BY-SA 4.0}.
The French original remains copyright of its authors and publishers.

\tableofcontents
\mainmatter
\chapter{\'Etale morphisms}
\label{I}

% original p. 1
To simplify the exposition, all the preschemes considered are assumed
to be locally noetherian, at least after no.~\ref{I.2}.

\section{Notions of differential calculus}
\label{I.1}

Let $X$ be a prescheme over~$Y$, and let $\Delta_{X/Y}$ or $\Delta$
be the diagonal morphism $X\to X\times_Y X$. It is an immersion,
hence a closed immersion of~$X$ into an open subset $V$ of~$X\times_Y X$.
Let $\cal{I}_X$ be the ideal of the closed subprescheme corresponding
to the diagonal in~$V$ (N.B. if one wants to proceed intrinsically,
without assuming $X$ separated over~$Y$---an assumption that would be
quite a joke---one should consider the set-theoretic inverse image
of~$\cal{O}_{X\times X}$ on~$X$, and denote by $\cal{I}_X$ the
augmentation ideal in the latter\ldots). The sheaf
$\cal{I}_X/\cal{I}_X^2$ can be regarded as a quasi-coherent sheaf
on~$X$; it is denoted by~$\mathit{\Omega}^1_{X/Y}$.
It is of finite type if $X\to Y$ is of finite type. It behaves well
with respect to extension of the base $Y'\to Y$. One also introduces
the sheaves $\cal{O}_{X\times_Y X}/\cal{I}_X^{n+1}=\cal{P}^n_{X/Y}$;
these are sheaves of \emph{rings} on~$X$, making $X$ into a prescheme
which one can denote by $\Delta_{X/Y}^n$ and call the
\emph{$n$th infinitesimal neighborhood of}~$X/Y$.
The corresponding sorites is utterly trivial, although rather
long\footnote{\cf EGA IV~16.3.}; it would be prudent to discuss it
only when one has something useful to say about it, in connection
with smooth morphisms.

\section{Quasi-finite morphisms}
\label{I.2}

\begin{proposition}
\label{I.2.1}
Let $A\to B$ be a local homomorphism (N.B. the rings are now
noetherian), and let $\goth{m}$ be the maximal ideal of~$A$.
The following conditions are equivalent:
\begin{enumerate}
\item[(i)] $B/\goth{m}B$ is finite-dimensional over~$k=A/\goth{m}$.
\item[(ii)] $\goth{m}B$ is an ideal of definition, and
$B/\goth{r}(B)=\kres(B)$ is an extension of~$k=\kres(A)$.
\item[(iii)] The completion $\widehat{B}$ is finite over the
completion $\widehat{A}$ of~$A$.
\end{enumerate}
\end{proposition}

One then says
% original p. 2
that $B$ \emph{is quasi-finite} over~$A$. A morphism $f\colon X\to Y$
is said to be quasi-finite at~$x$ (or the $Y$-prescheme $f$ is said
to be quasi-finite at~$x$) if $\cal{O}_x$ is quasi-finite
over~$\cal{O}_{f(x)}$. This also amounts to saying that $x$ is
\emph{isolated in its fiber}~$f^{-1}(x)$. A morphism is said to be
quasi-finite if it is so at every point\footnote{In EGA~II~6.2.3,
$f$ is also assumed to be of finite type.}.

\begin{corollary}
\label{I.2.2}
If $A$ is complete, quasi-finite is equivalent to finite.
\end{corollary}

One could give the usual sorites \textup{(i), (ii), (iii), (iv), (v),}
for quasi-finite morphisms, but this does not seem indispensable here.

\section{Unramified or net morphisms}
\label{I.3}

\begin{proposition}
\label{I.3.1}
Let $f\colon X\to Y$ be a morphism of finite type, $x\in X$,
$y=f(x)$. The following conditions are equivalent:
\begin{enumerate}
\item[(i)] $\cal{O}_x/\goth{m}_y\cal{O}_x$ is a finite separable
extension of~$\kres(y)$.
\item[(ii)] $\mathit{\Omega}^1_{X/Y}$ vanishes at~$x$.
\item[(iii)] The diagonal morphism $\Delta_{X/Y}$ is an open
immersion in a neighborhood of~$x$.
\end{enumerate}
\end{proposition}

For the implication (i)$\To$(ii), Nakayama immediately reduces
one to the case where $Y=\Spec(k)$, $X=\Spec(k')$, where this is
well known and, moreover, trivial from the definition of separability;
(ii)$\To$(iii) follows from a pleasant and easy characterization
of open immersions, using Krull; (iii)$\To$(i) because one is again
reduced to the case where $Y=\Spec(k)$ and the diagonal morphism
is an open immersion everywhere. One must then prove that $X$ is
finite with ring separable over~$k$; for this one is reduced to
the case where $k$ is algebraically closed. But then every closed
point of~$X$ is isolated (being identical to the inverse image of
the diagonal under the morphism $X\to X\times_k X$ defined by~$x$),
whence $X$ is finite. One may then suppose that $X$ consists of
a single point, with ring~$A$, so $A\otimes_k A\to A$ is an
isomorphism, whence $A=k$, qed.

\begin{definition}
\label{I.3.2}
\begin{enumerate}
\item[a)] One then says that $f$ is \emph{net}, or
\emph{unramified}, at~$x$, or that $X$ is net, or unramified,
at~$x$ over~$Y$.
\item[b)] Let $A\to B$ be a local homomorphism. One says that it
is \emph{net}, or \emph{unramified}, or that $B$ is a \emph{net},
or \emph{unramified}, local algebra over~$A$, if $B/\goth{r}(A)B$
is a finite separable extension of~$A/\goth{r}(A)$, \ie if
$\goth{r}(A)B=\goth{r}(B)$ and $\kres(B)$ is a separable extension
of~$\kres(A)$\footnote{\Cf the second thoughts in
III~\SourceRef{III.1.2}{1.2}.}.
\end{enumerate}
\end{definition}

\begin{remarks}
The fact
% original p. 3
that $B$ is net over~$A$ can already be recognized from the
completions of~$A$ and~$B$. Net implies quasi-finite.
\end{remarks}

\begin{corollary}
\label{I.3.3}
The set of points where $f$ is net is open.
\end{corollary}

\begin{corollary}
\label{I.3.4}
Let $X'$, $X$ be two preschemes of finite type over~$Y$, and
$g\colon X'\to X$ a $Y$-morphism. If $X$ is net over~$Y$, the
graph morphism $\Gamma_g\colon X'\to X\times_Y X$ is an open
immersion.
\end{corollary}

Indeed, it is the inverse image of the diagonal morphism
$X\to X\times_Y X$ under
\[
g\times_Y\id_{X'}\colon X'\times_Y X\to X\times_Y X.
\]

One can also introduce the annihilator ideal $\goth{d}_{X/Y}$
of~$\mathit{\Omega}^1_{X/Y}$, called the different ideal
of~$X/Y$; it defines a closed subprescheme of~$X$ whose underlying
set is the set of points where $X/Y$ is ramified, \ie not net.

\begin{proposition}
\label{I.3.5}
\begin{enumerate}
\item[(i)] An immersion is net.
\item[(ii)] The composite of two net morphisms is net.
\item[(iii)] Extension of the base in a net morphism gives another
net morphism.
\end{enumerate}
\end{proposition}

This can be seen equally well from (ii) or (iii) (the latter seems
more amusing to me). One can of course also be more precise by
giving pointwise statements; this is more general only in appearance
(except in the setting of definition~b)), and tedious. As usual,
one obtains corollaries:

\begin{corollaries}
\label{I.3.6}
\begin{enumerate}
\item[(iv)] A cartesian product of two net morphisms is another
net morphism.
\item[(v)] If $gf$ is net, $f$ is net.
\item[(vi)] If $f$ is net, $f_{\text{\textup{red}}}$ is net.
\end{enumerate}
\end{corollaries}

\begin{proposition}
\label{I.3.7}
Let $A\to B$ be a local homomorphism; suppose that the residue
field extension $\kres(B)/\kres(A)$ is trivial or that $\kres(A)$
is algebraically closed. For $B/A$ to be net, it is necessary and
sufficient that $\widehat{B}$ be (as an $\widehat{A}$-algebra)
a quotient of~$\widehat{A}$.
\end{proposition}

\begin{remarks}
\begin{enumerate}
\item[--] When the residue field extension is not assumed trivial,
one can reduce to this case by making a suitable finite flat
extension of~$A$ which kills the said extension.
\item[--] Give the example where $A$ is the local ring of an
ordinary double point of a curve, and $B$ that of a point of its
normalization: then $A\subset B$, $B$ is net over
% original p. 4
$A$ with trivial residue field extension, and
$\widehat{A}\to\widehat{B}$ is surjective but \emph{not injective}.
We shall therefore strengthen the notion of being net.
\end{enumerate}
\end{remarks}

\section{\'Etale morphisms. \'Etale coverings}
\label{I.4}

We shall admit everything we shall need concerning flat morphisms;
these facts will be proved later, if necessary\footnote{\Cf
Exp\ptbl\SourceRef{IV}{IV}.}.

\begin{definition}
\label{I.4.1}
\begin{enumerate}
\item[a)] Let $f\colon X\to Y$ be a morphism of finite type.
One says that $f$ is \emph{\'etale} at~$x$ if $f$ is flat at~$x$
and net at~$x$. One says that $f$ is \'etale if it is so at every
point. One says that $X$ is \'etale at~$x$ over~$Y$, or that it
is a $Y$-prescheme \'etale at~$x$, etc.
\item[b)] Let $f\colon A\to B$ be a local homomorphism. One says
that $f$ is \'etale, or that $B$ is \'etale over~$A$, if $B$ is
flat and unramified over~$A$.\footnote{\Cf the second thoughts in
III~\SourceRef{III.1.2}{1.2}.}
\end{enumerate}
\end{definition}

\begin{proposition}
\label{I.4.2}
For $B/A$ to be \'etale, it is necessary and sufficient that
$\widehat{B}/\widehat{A}$ be so.
\end{proposition}

In
% original p. 5
fact, this is true separately for ``net'' and for ``flat''.

\begin{corollary}
\label{I.4.3}
Let $f\colon X\to Y$ be of finite type, and $x\in X$. Whether
$f$ is \'etale at~$x$ depends only on the local homomorphism
$\cal{O}_{f(x)}\to\cal{O}_x$, and indeed only on the corresponding
homomorphism of completions.
\end{corollary}

\begin{corollary}
\label{I.4.4}
Suppose that the residue field extension $\kres(A)\to\kres(B)$
is trivial, or that $\kres(A)$ is algebraically closed. Then
$B/A$ is \'etale if and only if $\widehat{A}\to\widehat{B}$
is an isomorphism.
\end{corollary}

Combine flatness and~\ref{I.3.7}.

\begin{proposition}
\label{I.4.5}
Let $f\colon X\to Y$ be a morphism of finite type. Then the set
of points where it is \'etale is open.
\end{proposition}

Indeed, this is true separately for ``net'' and ``flat''.

This proposition shows that one can dispense with ``pointwise''
statements in studying morphisms of finite type which are \'etale
somewhere.

\begin{proposition}
\label{I.4.6}
\begin{enumerate}
\item[(i)] An open immersion is \'etale.
\item[(ii)] The composite of two \'etale morphisms is \'etale.
\item[(iii)] Extension of the base.
\end{enumerate}
\end{proposition}

Indeed, (i) is trivial, and for (ii) and (iii) it suffices to note
that this is true for ``net'' and for ``flat''. To tell the truth,
there are also corresponding statements for local homomorphisms
(without a finiteness condition), which in any case will have to
appear in the multiplodoque (beginning with the case: net).

\begin{corollary}
\label{I.4.7}
A cartesian product of two \'etale morphisms is likewise \'etale.
\end{corollary}

\begin{corollary}
\label{I.4.8}
Let $X$ and $X'$ be of finite type over~$Y$, and $g\colon X\to X'$
a $Y$-morphism. If $X'$ is unramified over~$Y$ and $X$ is \'etale
over~$Y$, then $g$ is \'etale.
\end{corollary}

Indeed, $g$ is the composite of the graph morphism
$\Gamma_g\colon X\to X\times_Y X'$, which is an open immersion
by~\ref{I.3.4}, and the projection morphism, which is \'etale
because it is obtained from the \'etale morphism $X\to Y$ by the
``change of base'' $X'\to Y$.

\begin{definition}
\label{I.4.9}
An \'etale (\resp net) covering of~$Y$ is a $Y$-scheme $X$ which
is finite over~$Y$ and \'etale (\resp net) over~$Y$.
\end{definition}

The first condition means that $X$ is defined by a coherent sheaf
of algebras $\cal{B}$ on~$Y$. The second then means that $\cal{B}$
is locally free on~$Y$ (\resp nothing at all), \emph{and} that,
in addition, for every $y\in Y$, the fiber
$\cal{B}(y)=\cal{B}_y\otimes_{\cal{O}_y}\kres(y)$ is a separable
algebra (= finite product of finite separable extensions)
over~$\kres(y)$.

\begin{proposition}
\label{I.4.10}
Let $X$ be a flat covering of~$Y$ of degree~$n$ (the definition
of this term deserved to appear in~\ref{I.4.9}), defined by a
coherent locally free sheaf $\cal{B}$ of algebras. One defines
in the well-known way the trace homomorphism $\cal{B}\to\cal{A}$
(which is a homomorphism of $\cal{A}$-modules, where
$\cal{A}=\cal{O}_Y$). For $X$ to be \'etale, it is necessary
and sufficient that the corresponding bilinear form
$\trace_{\cal{B}/\cal{A}}xy$ define an isomorphism of~$\cal{B}$
onto~$\cal{B}$, or, what amounts to the same, that the
\emph{discriminant section}
\[
d_{X/Y} = d_{\cal{B}/\cal{A}} \in
\Gamma(Y,\tbigwedge^n\check{\cal{B}}\otimes_{\cal{A}}
\tbigwedge^n\check{\cal{B}})
\]
be invertible, or finally that the discriminant ideal defined by
this section be the unit ideal.
\end{proposition}

Indeed, one is reduced to the case where $Y=\Spec(k)$, and then
this is a well-known separability criterion, trivial on passing
to the algebraic closure of~$k$.

\begin{remark}
We
% original p. 6
shall have a less trivial statement below, when one does not
assume a priori that $X$ is flat over~$Y$ but makes a normality
assumption.
\end{remark}

\section{The fundamental property of \'etale morphisms}
\label{I.5}

\begin{theorem}
\label{I.5.1}
Let $f\colon X\to Y$ be a morphism of finite type. For $f$ to be
an open immersion, it is necessary and sufficient that it be
an \emph{\'etale and radicial} morphism.
\end{theorem}

Recall that radicial means: injective, with purely inseparable
residue field extensions (and one may also recall that this means
that the morphism remains injective under every extension of the
base). Necessity is trivial; sufficiency remains. We shall give
two different proofs, the first shorter, the second more elementary.

1) A flat morphism is open, so one may suppose (replacing $Y$
by~$f(X)$) that $f$ is a \emph{homeomorphism} onto~$Y$.
Under any extension of the base, $f$ will remain flat, radicial,
and surjective, hence a homeomorphism, a fortiori closed.
Thus $f$ is \emph{proper}. Thus $f$ is \emph{finite}
(reference: Chevalley's theorem), defined by a coherent sheaf
$\mathcal{B}$ of algebras. $\mathcal{B}$ is locally free;
moreover, by the hypothesis it has rank~$1$ everywhere, so
$X=Y$, qed.

2) One may suppose that $Y$ and $X$ are \emph{affine}. One is
further easily reduced to proving the following: if $Y=\Spec(A)$,
with $A$ local, and if $f^{-1}(y)$ is nonempty ($y$ being the
closed point of~$Y$), then $X=Y$ (indeed, this will imply that
every $y\in f(X)$ has an open neighborhood $U$ such that $X|U=U$).
We shall have $X=\Spec(B)$; we want to prove that $A=B$. But for
this one is reduced to proving the analogous assertion on
replacing $A$ by~$\hat A$, and $B$ by $B\otimes_A\hat A$
(taking into account that $\hat A$ is faithfully flat over~$A$).
One may therefore suppose that $A$ is \emph{complete}. Let $x$
be the point above~$y$; by Corollary~\ref{I.2.2}, $\mathcal{O}_x$
is finite over~$A$ and hence (being flat and radicial over~$A$)
is identical to~$A$. Thus $X=Y\amalg X^\prime$ (disjoint sum).
Since $X$ is radicial over~$Y$, $X^\prime$ is empty. We are done.

\begin{corollary}
\label{I.5.2}
Let $f\colon X\to Y$ be a morphism which is a \emph{closed
immersion} and is \emph{\'etale}. If $X$ is connected, $f$ is
an isomorphism of~$X$ onto a connected component of~$Y$.
\end{corollary}

Indeed, $f$ is also an open immersion. One deduces:

\begin{corollary}
\label{I.5.3}
Let
% original p. 7
$X$ be a net $Y$-scheme, with $Y$ connected. Then every section
of~$X$ over~$Y$ is an isomorphism of~$Y$ onto a connected
component of~$X$. There is therefore a one-to-one correspondence
between the set of these sections and the set of connected
components $X_i$ of~$X$ such that the projection $X_i\to Y$
is an isomorphism (or, equivalently by~\ref{I.5.1}, is surjective
and radicial). In particular, a section is known when its value
at one point is known.
\end{corollary}

Only the first assertion requires a proof; by~\ref{I.5.2}, it
suffices to observe that a section is a closed immersion
(since $X$ is separated over~$Y$) and is \'etale by~\ref{I.4.8}.

\begin{corollary}
\label{I.5.4}
Let $X$ and $Y$ be two preschemes over~$S$, with $X$ net and
separated over~$S$ and $Y$ connected. Let $f$, $g$ be two
$S$-morphisms from~$Y$ to~$X$, and $y$ a point of~$Y$; suppose
that $f(y)=g(y)=x$ and that the residue field homomorphisms
$\kres(x)\to\kres(y)$ defined by $f$ and~$g$ are identical
(``$f$ and~$g$ agree geometrically at~$y$''). Then $f$ and~$g$
are identical.
\end{corollary}

This follows from~\ref{I.5.3} by reducing to the case where
$Y=S$, replacing $X$ by $X\times_S Y$.

Here is a particularly important variant of~\ref{I.5.3}.

\begin{theorem}
\label{I.5.5}
Let $S$ be a prescheme, $X$ and $Y$ two $S$-preschemes, $S_0$
a closed subprescheme of~$S$ with the same underlying space
as~$S$, and $X_0=X\times_S S_0$ and $Y_0=Y\times_S S_0$ the
``restrictions'' of~$X$ and~$Y$ to~$S_0$. Suppose that $X$ is
\'etale over~$S$. Then the natural map
\[
\Hom_S(Y,X)\to\Hom_{S_0}(Y_0,X_0)
\]
is \emph{bijective}.
\end{theorem}

One is again reduced to the case where $Y=S$, and then this
follows from the ``topological'' description of the sections
of~$X/Y$ given in~\ref{I.5.3}.

\begin{scholium}
This result contains an assertion of \emph{uniqueness} and
\emph{existence} of \emph{morphisms}. It can also be expressed
(when $X$ and~$Y$ are both taken \'etale over~$S$) by saying that
the functor $X\mto X_0$ from the category of \'etale $S$-schemes
to the category of \'etale $S_0$-schemes is \emph{fully faithful},
\ie establishes an \emph{equivalence} of the first with a
\emph{full subcategory} of the second. We shall see below that
it is even an equivalence of the first and the second
(which will be a theorem on the \emph{existence of \'etale
$S$-schemes}).
\end{scholium}

The
% original p. 8
following form of~\ref{I.5.5}, more general in appearance, is
often convenient:

\begin{corollary}[``Theorem on extending liftings'']
\label{I.5.6}
Consider a commutative diagram
\[
\xymatrix{
X\ar[d]&\ar[l]Y_0\ar[d]\\
S&\ar[l]Y
}
\]
of morphisms, where $X\to S$ is \'etale and $Y_0\to Y$ is
a bijective closed immersion. Then one can find a unique
morphism $Y\to X$ which makes the two corresponding triangles
commutative.
\end{corollary}

Indeed, replacing $S$ by~$Y$ and $X$ by $X\times_S Y$, one is
reduced to the case where $Y=S$, and then this is the special
case of~\ref{I.5.5} for $Y=S$.

Let us also point out the following immediate consequence
of~\ref{I.5.1} (which we did not give as Corollary~1 so as not
to interrupt the line of ideas developed after~\ref{I.5.1}):

\begin{proposition}
\label{I.5.7}
Let $X$, $X^\prime$ be two preschemes of finite type and flat
over~$Y$, and let $g\colon X\to X^\prime$ be a $Y$-morphism.
For $g$ to be an open immersion (\resp an isomorphism), it is
necessary and sufficient that, for every $y\in Y$, the induced
morphism on the fibers
\[
g\otimes_Y\kres(y)\colon X\otimes_Y\kres(y)\to X^\prime
\otimes_Y\kres(y)
\]
be so.
\end{proposition}

It suffices to prove sufficiency; since the assertion holds for
surjectivity, one is reduced to the case of an open immersion.
By~\ref{I.5.1}, one must check that $g$ is \emph{radicial},
which is trivial, and that it is \emph{\'etale}, which follows
from Corollary~\ref{I.5.9} below.

\begin{corollary}
\label{I.5.8}
(Should be moved to \No\ref{I.3}.) Let $X$ and $X^\prime$ be
two $Y$-preschemes, $g\colon X\to X^\prime$ a $Y$-morphism,
$x$ a point of~$X$, and $y$ its projection on~$Y$. For $g$ to
be quasi-finite (\resp net) at~$x$, it is necessary and sufficient
that $g\otimes_Y\kres(y)$ be so.
\end{corollary}

Indeed, the two algebras over~$k\big(g(x)\big)$ which one must
examine to check that one indeed has a morphism quasi-finite
\resp net at~$x$ are the same for~$g$ and $g\otimes_Y\kres(y)$.

\begin{corollary}
\label{I.5.9}
With
% original p. 9
the notation of~\ref{I.5.8}, suppose that $X$ and~$X^\prime$
are flat and of finite type over~$Y$. For $g$ to be flat
(\resp \'etale) at~$x$, it is necessary and sufficient that
$g\otimes_Y\kres(y)$ be so.
\end{corollary}

For ``flat'', the statement is included only as a reminder;
it is one of the fundamental flatness criteria\footnote{\Cf
IV~\SourceRef{IV.5.9}{5.9}.}. For \'etale, it follows from this,
taking~\ref{I.5.8} into account.

\section{Application to \'etale extensions of complete local rings}
\label{I.6}

This no.\ is a special case of results on formal preschemes
which will have to appear in the multiplodoque. Nevertheless,
one gets by here with less effort, \ie without the explicit
local determination of \'etale morphisms in
\No\ref{I.7} (using the Main Theorem). This is perhaps
sufficient reason to keep the present no.\ (even in the
multiplodoque) in this place.

\begin{theorem}
\label{I.6.1}
Let $A$ be a complete local ring (noetherian, of course), with
residue field~$k$. For every $A$-algebra~$B$, let
$R(B)=B\otimes_A k$, regarded as a $k$-algebra; it therefore
depends functorially on~$B$. Then $R$ defines an
\emph{equivalence} of the category of $A$-algebras
\emph{finite and \'etale over~$A$} with the category of
\emph{separable algebras of finite rank} over~$k$.
\end{theorem}

First of all, the functor in question is fully faithful, as
follows from the more general fact:

\begin{corollary}
\label{I.6.2}
Let $B$, $B^\prime$ be two $A$-algebras finite over~$A$. If $B$
is \'etale over~$A$, then the canonical map
\[
\Hom_{\textup{$A$-alg}}(B,B^\prime)\to
\Hom_{\textup{$k$-alg}}\big(R(B),R(B^\prime)\big)
\]
is bijective.
\end{corollary}

One is reduced to the case where $A$ is artinian (by replacing
$A$ by~$A/\goth{m}^n$), and then this is a special case
of~\ref{I.5.5}.

It remains to prove that for every finite separable $k$-algebra
(why not say: \'etale, it is shorter) $L$, there exists a $B$
\'etale over~$A$ such that $R(B)$ is isomorphic to~$L$. One may
suppose that $L$ is a separable extension of~$k$; as such, it
has a generator~$x$, \ie is isomorphic to an algebra $k[t]/Fk[t]$
where $F\in k[t]$ is a monic polynomial. One lifts~$F$ to a
monic polynomial $F_1$ in $A[t]$, and takes $B=A[t]/F_1A[t]$.

\section{Local construction of unramified and \'etale morphisms}
\label{I.7}
% original p. 10

\begin{proposition}
\label{I.7.1}
Let $A$ be a noetherian ring, $B$ a finite algebra over~$A$,
$u$ a generator of~$B$ over~$A$, $F\in A[t]$ such that $F(u)=0$
($F$ is not assumed monic), $u^\prime=F^\prime(u)$ (where
$F^\prime$ is the derivative polynomial), $\goth{q}$ a prime
ideal of~$B$ not containing~$u^\prime$, and $\goth{p}$ its
contraction to~$A$. Then $B_\goth{q}$ is net over~$A_\goth{p}$.
\end{proposition}

In other words, setting $Y=\Spec(A)$, $X=\Spec(B)$,
$X_{u^\prime}=\Spec(B_{u^\prime})$, $X_{u^\prime}$ is unramified
over~$Y$. The statement follows from the following, more precise one:

\begin{corollary}
\label{I.7.2}
The different ideal of~$B/A$ contains $u^\prime B$, and is equal
to it if the natural homomorphism $A[t]/FA[t]\to B$ (sending
$t$ to~$u$) is an isomorphism.
\end{corollary}

Let $J$ be the kernel of the homomorphism $C=A[t]\to B$; this
kernel contains $FA[t]$, and is equal to it in the second case
considered in~\ref{I.7.2}. Since the homomorphism is surjective,
$\Omega^1_{B/A}$ is identified with the quotient of~$\Omega^1_{C/A}$
by the submodule generated by $J\Omega^1_{C/A}$ and $d(J)$
(the definition of the homomorphism~$d$, and the calculation
of~$\Omega^1$ for a polynomial algebra, should have been made
explicit in \No\ref{I.1}). Identifying $\Omega^1_{C/A}$ with~$C$
by means of the basis~$dt$, one finds $B/B\cdot J^\prime$;
thus the different is generated by the set~$J^\prime$ of images
in~$B$ of the derivatives of the $G\in J$ (and it suffices to
take $G$ generating~$J$). Since $F\in J$, \resp $F$ is a generator
of~$J$, we are done. (N.B.\ One should make~\ref{I.7.2} a
proposition and~\ref{I.7.1} a corollary.) One obtains:

\begin{corollary}
\label{I.7.3}
Under the conditions of~\ref{I.7.1}, assuming that $F$ is monic
and that $A[t]/FA[t]\to B$ is an isomorphism, for $B_\goth{q}$
to be \'etale over~$A_\goth{p}$, it is necessary and sufficient
that~$\goth{q}$ not contain~$u^\prime$.
\end{corollary}

Indeed, since $B$ is flat over~$A$, \'etale is equivalent to
net, and one can apply~\ref{I.7.2}.

\begin{corollary}
\label{I.7.4}
Under the conditions of~\ref{I.7.3}, for $B$ to be \'etale
over~$A$, it is necessary and sufficient that~$u^\prime$ be
invertible, or equivalently that the ideal generated by $F$,
$F^\prime$ in $A[t]$ be the unit ideal.
\end{corollary}

The last criterion follows from the first and Nakayama (in~$B$).

A monic polynomial $F\in A[t]$ having the property stated in
Corollary~\ref{I.7.4}
% original p. 11
is called a \emph{separable polynomial} (if $F$ is not monic,
one should at least require its leading coefficient to be
invertible; when $A$ is a field, one recovers the usual definition).

\begin{corollary}
\label{I.7.5}
Let $B$ be a \emph{finite} algebra over the \emph{local} ring~$A$.
Suppose that $K(A)$ is infinite or that $B$ is local. Let~$n$
be the rank of~$L=B\otimes_A K(A)$ over~$K(A)=k$. For $B$ to
be net (\resp \'etale) over~$A$, it is necessary and sufficient
that $B$ be isomorphic to a quotient of (\resp isomorphic to)
$A[t]/FA[t]$, where $F$ is a monic separable polynomial which
one may assume to be (\resp which is necessarily) of degree~$n$.
\end{corollary}

Only necessity needs to be proved. Suppose that $B$ is net
over~$A$, so that $L$ is separable over~$k$. It then follows
from the hypothesis that $L/k$ has a generator~$\xi$, so the
$\xi^i$ ($0\leq i<n$) form a basis of~$L$ over~$k$. Let
$u\in B$ lift~$\xi$; then, by Nakayama, the $u^i$ ($0\leq i<n$)
generate the $A$-module~$B$ (\resp form a basis of it).
In particular, one can find a monic polynomial $F\in A[t]$
such that $F(u)=0$, and $B$ will be isomorphic to a quotient
of (\resp isomorphic to) $A[t]/FA[t]$. Finally, by~\ref{I.7.4}
applied to $L/k$, $F$ and~$F^\prime$ generate $A[t]$ modulo
$\goth{m}A[t]$; hence (by Nakayama in $A[t]/FA[t]$) $F$ and
$F^\prime$ generate~$A[t]$, and we are done.

\begin{theorem}
\label{I.7.6}
Let $A$ be a local ring, and $A\to\cal{O}$ a local homomorphism
such that~$\cal{O}$ is isomorphic to a localization of an algebra
of finite type over~$A$. Suppose that $\cal{O}$ is \emph{net}
over~$A$. Then one can find an $A$-algebra $B$, integral over~$A$,
a maximal ideal $\goth{n}$ of~$B$, a generator $u$ of~$B$
over~$A$, and a monic polynomial $F\in A[t]$, such that
$\goth{n}\not\ni F^\prime(u)$ and $\cal{O}$ is isomorphic
(as an $A$-algebra) to~$B_{\goth{n}}$. If $\cal{O}$ is \'etale
over~$A$, one can take $B=A[t]/FA[t]$.
\end{theorem}

(Of course, these conditions are also sufficient\ldots)

Let us first point out the pleasant corollaries:

\begin{corollary}
\label{I.7.7}
For $\cal{O}$ to be net over~$A$, it is necessary and sufficient
that $\cal{O}$ be isomorphic to the quotient of an analogous
algebra which is \emph{\'etale} over~$A$.
\end{corollary}

Indeed, one takes $\cal{O}^\prime=B^\prime_{\goth{n}^\prime}$,
where $B^\prime=A[t]/FA[t]$ and $\goth{n}^\prime$ is the inverse
image of~$\goth{n}$ in~$B^\prime$.

\begin{corollary}
\label{I.7.8}
Let
% original p. 12
$f\colon X\to Y$ be a morphism of finite type, $x\in X$.
For $f$ to be net at~$x$, it is necessary and sufficient that
there exist an open neighborhood $U$ of~$x$ such that $f|U$
factors as $U\to X^\prime\to Y$, where the first arrow is
a closed immersion and the second an \'etale morphism.
\end{corollary}

This is a mere translation of~\ref{I.7.7}.

Let us show how the jargon of~\ref{I.7.6} follows from the
main statement: indeed, by~\ref{I.7.7} there is an epimorphism
$\cal{O}^\prime\to\cal{O}$, where $\cal{O}$ has the desired
properties; but since $\cal{O}^\prime$ and~$\cal{O}$ are \'etale
over~$A$, the morphism $\cal{O}^\prime\to\cal{O}$ is \'etale
by~\ref{I.4.8}, hence an isomorphism.

\subsubsection*{Proof of~\ref{I.7.6}}
It follows a proof from the Chevalley seminar. By the
\emph{Main Theorem}, one has $\cal{O}=B_{\goth{n}}$, where
$B$ is a finite algebra over~$A$ and $\goth{n}$ is a maximal
ideal of it. Then $B/\goth{n}=K(\cal{O})$ is a separable,
hence monogenic, extension of~$k$; if $\goth{n}_i$
($1\leq i\leq r$) are the maximal ideals of~$B$ distinct
from~$\goth{n}$, there is therefore an element $u$ of~$B$
which belongs to all the~$\goth{n}_i$, and whose image in
$B/\goth{n}$ is a generator of it. Now
$B/\goth{n}=B_\goth{n}/\goth{n}B_\goth{n}
=B_\goth{n}/\goth{m}B_\goth{n}$ (where $\goth{m}$ is the
maximal ideal of~$A$). Let us admit for a moment the following

\begin{lemma}
\label{I.7.9}
Let $A$ be a local ring, $B$ a finite algebra over~$A$,
$\goth{n}$ a maximal ideal of~$B$, and $u$ an element of~$B$
whose image in $B_\goth{n}/\goth{m}B_\goth{n}$ generates it
as an algebra over~$k=A/\goth{m}$, and which belongs to every
maximal ideal of~$B$ distinct from~$\goth{n}$. Let
$B^\prime=B[u]$, $\goth{n}^\prime=\goth{n}B^\prime$.
Then the canonical homomorphism
$B^\prime_{\goth{n}^\prime}\to B_\goth{n}$ is an isomorphism.
\end{lemma}

\begin{lemma}
\label{I.7.10}
(Should have appeared as a corollary to~\ref{I.7.1},
before~\ref{I.7.5}, which it implies.) Let $B$ be a finite
algebra over~$A$ generated by an element~$u$, and let
$\goth{n}$ be a maximal ideal of~$B$ such that $B_\goth{n}$
is unramified over~$A$. Then there exists a monic polynomial
$F\in A[t]$ such that $F(u)=0$ and $F^\prime(u)\notin\goth{n}$.
\end{lemma}

Indeed, let $n$ be the rank of the $k$-algebra $L=B\otimes_A k$;
by Nakayama there exists a monic polynomial of degree~$n$
in $A[t]$ such that $F(u)=0$. Let $f$ be the polynomial obtained
from~$F$ by reduction mod~$\goth{m}$; then $L$ is $k$-isomorphic
to $k[t]/fk[t]$, so by~\ref{I.7.3}, $f^\prime(\xi)$ does not
belong to the maximal ideal of~$L$ corresponding to~$\goth{n}$
($\xi$ denoting the image of~$t$ in~$L$, \ie the image of~$u$
in~$L$). Since $f^\prime(\xi)$ is the image of~$F^\prime(u)$,
we are done.

The
% original p. 13
theorem~\ref{I.7.6} now follows by combining~\ref{I.7.9}
and~\ref{I.7.10}. It remains to prove~\ref{I.7.9}. Set
$S^\prime=B^\prime-\goth{n}^\prime$, so
$B^\prime{S^\prime}^{-1}=B^\prime_{\goth{n}^\prime}$.
\[
\begin{array}{ccccc}
B&\longrightarrow&B{S^\prime}^{-1}&\longrightarrow&BS^{-1} =
B_\goth{n}\\ \Big\uparrow&&\Big\uparrow&&\\
B^\prime&\longrightarrow&B^\prime{S^\prime}^{-1}
\rlap{\hbox{$\displaystyle \ = B^\prime_{\goth{n}^\prime}$}}&&
\end{array}
\]
Similarly, let $S=B-\goth{n}$, so $BS^{-1}=B_\goth{n}$.
There is therefore a natural homomorphism
$B{S^\prime}^{-1}\to BS^{-1}=B_\goth{n}$; let us prove that
it is an isomorphism, \ie that the elements of~$S$ are invertible
in $B{S^\prime}^{-1}$, \ie that every maximal ideal $\goth{p}$
of the latter is disjoint from~$S$, \ie contracts to~$\goth{n}$
in~$B$. Indeed, since $B{S^\prime}^{-1}$ is finite over
$B^\prime{S^\prime}^{-1}=B^\prime_{\goth{n}^\prime}$,
$\goth{p}$ contracts to the unique maximal ideal
$\goth{n}^\prime B_{\goth{n}^\prime}$ of
$B^\prime_{\goth{n}^\prime}$, hence contracts to the maximal
ideal $\goth{n}^\prime$ of~$B^\prime$; since $B$ is finite
over~$B^\prime$, the ideal $\goth{q}$ of~$B$ induced
by~$\goth{p}$, lying over~$\goth{n}^\prime$, is necessarily
maximal, and does not contain~$u$, hence is identical
to~$\goth{n}$. (We have just used the fact that $u$ belongs
to every maximal ideal of~$B$ distinct from~$\goth{n}$.)
Let us now prove that $B{S^\prime}^{-1}$ equals
$B^\prime{S^\prime}^{-1}$: since it is finite over the latter,
Nakayama reduces one to proving equality mod
$\goth{n}^\prime B{S^\prime}^{-1}$, and a fortiori it suffices
to prove equality mod $\goth{m}B{S^\prime}^{-1}$; now
$B{S^\prime}^{-1}/\goth{m}B{S^\prime}^{-1}
=B_\goth{n}/\goth{m}B_\goth{n}$ is generated over~$k$ by~$u$
(here one uses the other property of~$u$), so the image
of~$B^\prime$ (and a fortiori of $B^\prime{S^\prime}^{-1}$)
in it is the whole ring (being a subring containing~$k$
and the image of~$u$).

\begin{remark}
It should be possible to state Theorem~\ref{I.7.6} for a ring
$\cal{O}$ which is merely semilocal, so as to include~\ref{I.7.5}
as well: one would assume that $\cal{O}/\goth{m}\cal{O}$ is a
\emph{monogenic} $k$-algebra; one could therefore find a
$u\in B$ whose image in $B/\goth{m}B$ is a generator, and
which belongs to all the maximal ideals of~$B$ not coming
from~$\cal{O}$. Lemmas~\ref{I.7.9} and~\ref{I.7.10} should
adapt without difficulty. More generally, \ldots
\end{remark}

\section[Infinitesimal lifting of \'etale schemes]{Infinitesimal lifting of \'etale schemes.
Application to formal schemes}
\label{I.8}

\begin{proposition}
\label{I.8.1}
Let $Y$ be a prescheme, $Y_0$ a subprescheme, $X_0$ an \'etale
$Y_0$-scheme, and $x$ a point of~$X_0$. Then there exist an
\'etale $Y$-scheme~$X$, a neighborhood $U_0$ of~$x$ in~$X_0$,
and a $Y_0$-isomorphism $U_0\isomto X\times_Y Y_0$.
\end{proposition}

Indeed, let $y$ be the projection of~$x$ on~$Y_0$; applying
\ref{I.7.6} to the \'etale local homomorphism $A_0\to B_0$
of the local rings of~$y$ and~$x$ in $Y_0$ and~$X_0$, one
finds an isomorphism
% original p. 14
\[
B_0 = {C_0}_{\goth{n}_0}\qquad C_0 = A_0[t]/F_0A_0[t]
\]
where $F_0$ is a monic polynomial and $\goth{n}_0$ is a
maximal ideal of~$C_0$ not containing the class of $F^\prime_0(t)$
in~$C_0$. Let $A$ be the local ring of~$y$ in~$Y$, let $F$
be a monic polynomial in $A[t]$ giving~$F_0$ under the
surjective homomorphism $A\to A_0$ (one lifts the coefficients
of~$F_0$), and finally let $C=A[t]/FA[t]$ and $\goth{n}$ be
the maximal ideal of~$C$ which is the inverse image of
$\goth{n}_0$ under the natural epimorphism
$C\to C\otimes_A A_0=C_0$. Set
\[
B = C_\goth{n}.
\]
It is immediate from the construction and~\ref{I.7.1} that
$B$ is \'etale over~$A$, and that there is an isomorphism
$B\otimes_A A_0=A_0$. One knows (Chap\ptbl I) that there
exist a $Y$-scheme~$X$ of finite type and a point~$z$ of~$X$
above~$y$ such that $\cal{O}_z$ is $A$-isomorphic to~$C$;
since the latter is \'etale over~$A=\cal{O}_y$, one may
(taking $X$ small enough) suppose that $X$ is \'etale over~$Y$.
Let $X_0^\prime=X\times_Y Y_0$; then the local ring of~$z$
in~$X_0^\prime$ is identified with
$\cal{O}_z\otimes_A A_0=B\otimes_A A_0$, hence is isomorphic
to~$B_0$. This isomorphism is defined by an isomorphism of a
neighborhood~$U_0$ of~$x$ in~$X_0$ onto a neighborhood of~$z$
in~$X_0^\prime$ (loc.\ cit.), which one may suppose identical
to~$X_0^\prime$ by taking $X$ small enough. We are done.

\begin{corollary}
\label{I.8.2}
An analogous statement holds for \'etale \emph{coverings},
assuming that the residue field $\kres(y)$ is infinite.
\end{corollary}

The proof is the same, with~\ref{I.7.5} replacing~\ref{I.7.6}.

\begin{theorem}
\label{I.8.3}
The functor considered in~\ref{I.5.5} is an \emph{equivalence}
of \emph{categories}.
\end{theorem}

By Theorem~\ref{I.5.5}, it remains to show that every \'etale
$S_0$-scheme~$X_0$ is isomorphic to an $S_0$-scheme
$X\times_S S_0$, where $X$ is an \'etale $S$-scheme.
The topological space underlying~$X$ must necessarily be
identical to that of~$X_0$, with $X_0$ moreover identified
with a closed subprescheme of~$X$. The problem is therefore
equivalent to the following: to find, on the topological
space $|X_0|$ underlying~$X_0$, a sheaf of algebras~$\cal{O}_X$
over~$f_0^*(\cal{O}_S)$ (where $f_0$ is the projection
$X_0\to S_0$, regarded here as a continuous map of the
underlying spaces), making $|X_0|$ into an \'etale
$S$-prescheme~$X$, and a homomorphism of algebras
$\cal{O}_X\to\cal{O}_{X_0}$, compatible with the homomorphism
$f_0^*(\cal{O}_S)\to f_0^*(\cal{O}_{S_0})$ of sheaves of
scalars, inducing an isomorphism
$\cal{O}_X\otimes_{f_0^*(\cal{O}_S)}f_0^*(\cal{O}_{S_0})
\isomto\cal{O}_{X_0}$. (Then $X$ will be an \'etale
$S$-prescheme reducing to~$X_0$,
% original p. 15
hence will be separated over~$S$ since $X_0$ is separated
over~$S_0$, and $X$ answers the question.) Moreover, if
$(U_i)$ is a covering of~$X_0$ by open subsets, and if one
has found a solution of the problem in each~$U_i$, it follows
from the uniqueness theorem~\ref{I.5.5} that these solutions
glue (\ie the sheaves of algebras defining them, equipped
with their augmentation homomorphisms, glue), and one sees
at once that the ringed space thus constructed over~$S$ is
an \'etale $S$-prescheme~$X$ equipped with an isomorphism
$X\times_S S_0\isomfrom X_0$. It therefore suffices to find
a solution locally, which is ensured by~\ref{I.8.1}.

\begin{corollary}
\label{I.8.4}
Let $S$ be a locally noetherian formal prescheme, equipped
with an ideal of definition~$J$, and let
$S_0=\big(|S|,\cal{O}_S/\cal{J}\big)$ be the corresponding
ordinary prescheme. Then the functor
$\goth{X}\mto\goth{X}\times_S S_0$ from the category of
\'etale coverings of~$S$ to the category of \'etale coverings
of~$S_0$ is an equivalence of categories.
\end{corollary}

Of course, by an \'etale covering of a \emph{formal}
prescheme~$S$ one means a covering of~$S$, \ie a formal
prescheme over~$S$ defined by means of a coherent sheaf
of algebras~$\cal{B}$, such that $\cal{B}$ is \emph{locally
free} and the residue fibers
$\cal{B}_s\otimes_{\cal{O}_s}\kres(s)$ of~$\cal{B}$ are
\emph{separable} algebras over~$\kres(s)$. If one denotes
by~$S_n$ the ordinary prescheme
$\big(|S|,\cal{O}_S/J^{n+1}\big)$, the datum of a coherent
sheaf of algebras $\cal{B}$ on~$S$ is equivalent to the datum
of a sequence of coherent sheaves of algebras $\cal{B}_n$
on the~$S_n$, equipped with a transitive system of
homomorphisms $\cal{B}_m\to\cal{B}_n$ ($m\geq n$) defining
isomorphisms
$\cal{B}_m\otimes_{\cal{O}_{S_m}}\cal{O}_{S_n}\isomto\cal{B}_n$.
It is immediate that $\cal{B}$ is locally free if and only
if the $\cal{B}_n$ on the~$S_n$ are so, and that the
separability condition holds if and only if it holds
for~$\cal{B}_0$, or equivalently for all the~$\cal{B}_n$.
Thus $\cal{B}$ is \'etale over~$S$ if and only if the
$\cal{B}_n$ on the~$S_n$ are so. In view of this,
\ref{I.8.4} follows at once from~\ref{I.8.3}.

\begin{remark}
It was not necessary in~\ref{I.8.4} to restrict oneself
to \emph{coverings}. This is, however, the only case used
for the moment.
\end{remark}

\section{Permanence properties}
\label{I.9}

Let $A\to B$ be a local \'etale homomorphism. We examine here
some cases in which a certain property of $A$ implies the
same property for~$B$, or conversely.
% original p. 16
A number of such propositions already follow merely from
the fact that $B$ is \emph{quasi-finite} and \emph{flat}
over~$A$, and we shall confine ourselves to ``recalling''
a few of them: $A$ \emph{and} $B$ \emph{have the same Krull
dimension and the same depth} (Serre's ``cohomological codimension'',
in the terminology still current). It follows, for example,
that \emph{$A$ is Cohen--Macaulay if and only if $B$ is so}.
Moreover, for every prime ideal $\goth{q}$ of~$B$, inducing
$\goth{p}$ on~$A$, $B_{\goth{q}}$ will again be quasi-finite
and flat over~$A_{\goth{p}}$, provided that $B$ is assumed
to be a localization of an algebra of finite type over~$A$
(this follows from the fact that the set of points where
a morphism of finite type is quasi-finite \resp flat is
open); furthermore, \emph{every} prime ideal $\goth{p}$ of~$A$
is induced by a prime ideal $\goth{q}$ of~$B$ (since $B$ is
\emph{faithfully} flat over~$A$). It follows, for example,
that \emph{$\goth{p}$ and~$\goth{q}$ have the same height};
and also that \emph{$A$ has no embedded prime ideals if
and only if~$B$ has none}.

We shall therefore restrict ourselves to propositions more
specific to the case of \'etale morphisms.

\begin{proposition}
\label{I.9.1}
Let $A\to B$ be a local \'etale homomorphism. For $A$ to be
regular, it is necessary and sufficient that $B$ be so.
\end{proposition}

Indeed, let $k$ be the residue field of~$A$, and $L$ that
of~$B$. Since $B$ is flat over~$A$ and $L=B\otimes_A k$,
\ie $\goth{n}=\goth{m}B$ (where $\goth{m}$,~$\goth{n}$
are the maximal ideals of~$A$,~$B$), the $\goth{m}$-adic
filtration on~$B$ is identical to its $\goth{n}$-adic
filtration, and one has
\[
\gr^*(B)=\gr^*(A)\otimes_k L.
\]
It follows that $\gr^*(B)$ is a polynomial algebra over~$L$
if and only if $\gr^*(A)$ is a polynomial algebra over~$k$,
qed. (N.B.\ The fact that $L/k$ is separable has not been used.)

\begin{corollary}
\label{I.9.2}
Let $f\colon X\to Y$ be an \'etale morphism. If $Y$ is
regular, so is $X$; the converse holds if $f$ is surjective.
\end{corollary}

% The corrected source repeats 9.2. Preserve it with a distinct PDF anchor.
\setcounter{theorem}{1}
\begingroup
\renewcommand{\theHtheorem}{\theHsection.\arabic{theorem}.bis}
\begin{proposition}
\label{prop:I.9.2}
Let $f\colon X\to Y$ be an \'etale morphism. If $Y$ is
reduced, so is~$X$; the converse holds if $f$ is surjective.
\end{proposition}
\endgroup

This is equivalent to the following

\begin{corollary}
\label{I.9.3}
Let $f\colon A\to B$ be a local \'etale homomorphism, with
$B$ isomorphic to an
% original p. 17
$A$-algebra which is a localization of an $A$-algebra of
finite type. For $A$ to be reduced, it is necessary and
sufficient that $B$ be so.
\end{corollary}

Necessity is trivial, since $A\to B$ is injective ($B$ being
faithfully flat over~$A$). Sufficiency: let $\goth{p}_i$ be
the minimal prime ideals of~$A$; by hypothesis the natural
map $A\to\prod_i A/\goth{p}_i$ is injective, so tensoring
with the flat $A$-module $B$, one finds that
$B\to\prod_i B/\goth{p}_iB$ is injective, and one is reduced
to proving that the $B/\goth{p}_iB$ are reduced. Since
$B/\goth{p}_iB$ is \'etale over~$A/\goth{p}_i$, one is
reduced to the case where $A$ is an integral domain. Let~$K$
be its field of fractions; since $A\to K$ is injective,
so is $B\to B\otimes_A K$ ($B$ being $A$-flat), and one is
reduced to proving that the latter ring is reduced. Now,
$B$, being a localization of an $A$-algebra of finite type
over~$A$, is the local ring of a point~$x$ of a scheme
$X=\Spec(C)$ of finite type and \emph{\'etale} over
$Y=\Spec(A)$; hence $B\otimes_A K$ is a localization (with
respect to a suitable multiplicatively closed set) of the
ring $C\otimes_A K$ of~$X\otimes_A K$. Since $X\otimes_A K$
is \'etale over~$K$, its ring is a finite product of fields
(separable extensions of~$K$); hence the same holds
for~$B\otimes_A K$, qed.

\begin{corollary}
\label{I.9.4}
Let $f\colon A\to B$ be a local \'etale homomorphism, and
suppose that $A$ is analytically reduced (\ie the completion
$\hat{A}$ of~$A$ has no nilpotent elements). Then $B$ is
analytically reduced, and a fortiori reduced.
\end{corollary}

Indeed, $\hat{B}$ is \emph{finite} and \'etale over~$\hat{A}$,
and one applies~\ref{I.9.3}.

\begin{theorem}
\label{I.9.5}
Let $f\colon A\to B$ be a local homomorphism, with $B$
isomorphic to a localization of an $A$-algebra of finite
type. Then:
\begin{enumerate}
\item[(i)] If $f$ is \'etale, $A$ is normal if and only if
$B$ is so.
\item[(ii)] If $A$ is normal, $f$ is \'etale if and only if
$f$ is injective and net (and then $B$ is normal by \textup{(i)}).
\end{enumerate}
\end{theorem}

We shall give two different proofs of (i). The first uses
some of the properties of flat quasi-finite morphisms
(recalled at the beginning of this no.) without using
\ref{I.7.6} (and thereby the Main Theorem); the reverse
is true of the second proof. Finally, for~(ii) it seems
that one needs the Main Theorem in any case.

\subsubsection*{First proof}
One uses the following necessary and sufficient condition
% original p. 18
for normality of a noetherian local ring $A$ of dimension
$\neq0$.

\begin{namedstatement}{Serre's criterion}
\textup{(i)}~For every prime ideal $\goth{p}$ of~$A$ of
height~$1$, $A_{\goth{p}}$ is normal (or, equivalently,
regular); \textup{(ii)}~For every prime ideal $\goth{p}$
of~$A$ of height~$\geq2$, the depth of $A_{\goth{p}}$ is
$\geq2$.\footnote{\Cf EGA~IV~5.8.6.}
\end{namedstatement}

We shall admit this criterion here; it is supposed to appear
in the paragraph on flat morphisms. Its main advantage
is that it does not assume a priori that $A$ is reduced,
let alone an integral domain. Here one may already suppose
that $\dim{A}=\dim{B}\neq0$.

By the facts recalled at the beginning of this no., the
prime ideals $\goth{p}$ of~$A$ of height~$1$ (\resp of
height~$\ge2$) are precisely the contractions to~$A$ of
the prime ideals~$\goth{q}$ of~$B$ of height~$1$ (\resp
of height~$\ge2$). Finally, if $\goth{p}$ and $\goth{q}$
correspond, $B_{\goth{q}}$ is \'etale over~$A_{\goth{p}}$,
hence has the same depth as $A_{\goth{p}}$, and is regular
if and only if $A_{\goth{p}}$ is so (\ref{I.9.1}). Applying
Serre's criterion, one finds that $A$ is normal if and only
if $B$ is so.

\subsubsection*{Second proof}
Suppose that $B$ is normal; let $L$ be its field of fractions,
and $K$ that of $A$ ($A$ is an integral domain since $B$ is
so). We saw in the proof of~\ref{I.9.3} that $B\otimes_A K$
is a finite product of fields; since it is contained in $L$
it is a field, and since it contains $B$ it is $L$. An
element of~$K$ integral over~$A$ is integral over~$B$, hence
belongs to $B$ since $B$ is normal, and therefore to $A$
since $B\cap K=A$ (as follows from the fact that $B$ is
faithfully flat over~$A$).

Suppose now that $A$ is normal; let us prove that $B$ is
so. By~\ref{I.7.6}, one has $B=B^\prime_{\goth{n}}$, where
$B^\prime=A[t]/FA[t]$, with $F$ and $\goth{n}$ as
in~\ref{I.7.6}. Thus $L=B\otimes_A K$ will be a localization
of~$B^\prime\otimes_A K=K[t]/FK[t]$, and a product of fields,
finite separable extensions of~$K$; this product (as whenever
one localizes an artinian ring, here $B^\prime_K$, with
respect to a multiplicatively closed set) is a direct
factor of~$B^\prime_K$, thus corresponding to a decomposition
$F=F_1F_2$ in $K[t]$, with the generator of~$L$ corresponding
to $t$ already annihilated by~$F_1$. Now, \emph{since $A$
is normal, the $F_i$ belong to $A[t]$} (assuming them monic).
Observing that $B\to L=B\otimes_A K$ is injective (since
$A\to K$ is so and $B$ is flat over~$A$), it follows that
one already has $F_1(u)=0$, with $u$ the class of~$t$ in~$L$.
Assuming that $F$ was chosen of minimal degree, it follows
that $F_2=1$ (N.B.\ One has
$F^\prime(u)=F^\prime_1(u)F_2(u)+F_1(u)F^\prime_2(u)
=F^\prime_1(u)F_2(u)$ since $F_1(u)=0$, whence
$F^\prime_1(u)\neq0$ since $F^\prime(u)\neq0$.)

Thus
% original p. 19
one has
\begin{equation*}
\label{eq:I.9.5.*}
\tag{$*$} L=B\otimes_A K=K[t]/FK[t]
\end{equation*}
and $F$ is consequently a separable polynomial in $K[t]$
(but of course not necessarily in $A[t]$). (N.B.\ For the
moment, we have essentially shown only that in~\ref{I.7.6}
one can choose $F$ and $\goth{n}$ in such a way that---with
the notation used here---$B^\prime\to B^\prime_{\goth{n}}=B$
is \emph{injective}; we used the normality of $A$ for this;
I do not know whether it remains true without the normality
assumption.)

Let us now recall the well-known lemma taken from Serre's
course of last year:

\begin{lemma}
\label{I.9.6}
Let $K$ be a ring, $F\in K[t]$ a monic separable polynomial,
$L=K[t]/FK[t]$, and $u$ the class of~$t$ in~$L$ (so that
$F^\prime(u)$ is an invertible element of~$L$). Then one
has the formulas (where $n=\deg{F}$):
\begin{enumerate}
\item[] $\trace_{L/K}u^i/F^\prime(u)=0$\quad if\quad $0\le i<n-1$,
\item[] $\trace_{L/K}u^{n-1}/F^\prime(u)=1$.
\end{enumerate}
\end{lemma}

\begin{corollary}
\label{I.9.7}
The determinant of the matrix
$(u^j\cdot u^i/F^\prime(u))_{0\le i,j\le n-1}$ is equal
to $(-1)^{n(n-1)/2}$, hence is invertible in every subring
$A$ of~$K$.
\end{corollary}

\begin{corollary}
\label{I.9.8}
Let $A$ be a subring of~$K$, $V$ the $A$-module generated
by the $u^i$ ($0\le i\le n-1$) in $L$, and $V^\prime$
the $A$-submodule of~$L$ consisting of the $x\in L$ such
that $\trace_{L/K}(xy)\in A$ for every $y\in V$ (\ie for
$y$ of the form $u^i$, $0\le i\le n-1$). Then $V^\prime$
is the $A$-module with basis the $u^i/F^\prime(u)$
($0\le i\le n-1$).
\end{corollary}

\begin{corollary}
\label{I.9.9}
Suppose that $K$ is the field of fractions of a normal
integral domain $A$, with $F$ having its coefficients
in $A$. Then, with the notation of~\ref{I.9.8}, $V^\prime$
contains the integral closure $A^\prime$ of~$A$ in $L$,
which is therefore contained in $A[u]/F^\prime(u)$ and
a fortiori in $A[u][F^\prime(u)^{-1}]$.
\end{corollary}

Let us apply this last corollary to the situation obtained
in the proof: since $F^\prime(u)$ is invertible in $B$,
which contains $A[u]$, $B$ contains $A^\prime$. By the
Main Theorem (or from the fact that $B=A[u]_{\goth{n}}$),
$B$ is a localization of $A^\prime$. Since $A^\prime$ is
normal, so is~$B$.

\subsubsection*{Proof of \textup{(ii)}}
One
% original p. 20
proceeds as in the preceding proof to show that, in
\ref{I.7.6}, one can choose $F$ so that ($*$) still holds.
The only obstacle a priori is that, since $B$ is no longer
assumed flat over~$A$, one can no longer assert that
$B\to L$ is injective, so that the argument applies a
priori only to the image $B_1$ of~$B$ under the said
homomorphism. It follows at once that $B_1$ is flat
over~$A$ (being a localization of a free algebra over~$A$).
By~\ref{I.4.8}, the morphism $B\to B_1$ is \'etale,
hence an isomorphism, which completes the proof.

(From the point of view of exposition, one should
interchange the last two proofs and put the formal
calculations of the lemma and its corollaries in a
separate no.)

\begin{corollary}
\label{I.9.10}
Let $f\colon X\to Y$ be an \'etale morphism. If $Y$ is
normal, so is $X$; the converse holds if $f$ is surjective.
\end{corollary}

\begin{corollary}
\label{I.9.11}
Let $f\colon X\to Y$ be a dominant morphism, with $Y$
normal and $X$ connected. If $f$ is net, $f$ is \'etale,
hence $X$ is normal and consequently (being connected)
irreducible.
\end{corollary}

Let $U$ be the set of points where $f$ is \'etale; it is
open, and it suffices to show that it is also closed and
nonempty. $U$ contains the inverse image of the generic
point of~$Y$ (since, for an algebra over a field,
unramified $=$ \'etale), hence ($X$ dominating $Y$) is
nonempty. If $x$ belongs to the closure of~$U$, then it
belongs to the closure of an irreducible component $U_i$
of~$U$, hence to an irreducible component
$X_i\overset{v}{=}\bar{U}_i$ of~$X$ which meets $U$ and
consequently dominates $Y$ (since every component of~$U$,
flat over~$Y$, dominates $Y$). Thus, if $y$ is the
projection of~$x$ on~$Y$, $\cal{O}_y\to\cal{O}_x$ is
\emph{injective} (taking into account that $\cal{O}_y$
is an integral domain). Since $\cal{O}_y$ is normal
and $\cal{O}_y\to\cal{O}_x$ is net, one concludes
using~\ref{I.9.5}(ii).

\begin{corollary}
\label{I.9.12}
Let $f\colon X\to Y$ be a dominant morphism of finite
type, with $Y$ normal and $X$ irreducible. Then the set
of points where $f$ is \'etale is identical to the
complement of the support of $\Omega^1_{X/Y}$, \ie to
the complement of the subprescheme of~$X$ defined by
the different ideal $\goth{d}_{X/Y}$.
\end{corollary}

(This
% original p. 21
is the ``less trivial'' statement alluded to in the
remark in~\No\ref{I.4}.)

\begin{remark}
One should beware of believing that a connected \'etale
covering of an irreducible scheme is itself irreducible
when the base is not assumed normal. This question
will be studied in~\No\ref{I.11}.
\end{remark}

\section{\'Etale coverings of a normal scheme}
\label{I.10}

\begin{proposition}
\label{I.10.1}
Let $X$ be an \'etale separated prescheme over a normal
connected $Y$ with field~$K$. Then the
connected components $X_i$ of~$X$ are integral, their
fields $K_i$ are finite separable extensions of~$K$,
and $X_i$ is identified with a nonempty open subset of
the normalization of~$X$ in $K_i$ (hence $X$ with a
dense open subset of the normalization of~$Y$ in
$R(X)=L=\prod K_i$).
\end{proposition}

By~\ref{I.9.10}, $X$ is normal; a fortiori its local
rings are integral domains, so the connected components
of~$X$ are irreducible. Since $X_i$ is normal, and finite
and dominant over~$Y$, it follows from a special case
(moreover, almost trivial) of the Main Theorem that
$X_i$ is an open subset of the normalization of~$X$
in the field $K_i$ of~$X$.

\begin{corollary}
\label{I.10.2}
Under the conditions of~\ref{I.10.1}, $X$ is finite
over~$Y$ (\ie an \'etale covering of~$Y$) if and only
if $X$ is isomorphic to the normalization $Y'$ of~$Y$
in $L=R(X)$ (the ring of rational functions on~$X$).
\end{corollary}

Indeed, one knows that this normalization is finite
over~$Y$ ($Y$ being normal and $R/K$ separable);
conversely, if $X$ is finite over~$Y$, it is finite
over~$Y'$, so its image in $Y'$ is closed, and on
the other hand it is dense.

An algebra $L$ of finite rank over~$K$ will be called
\emph{unramified over~$X$} (or simply unramified
over~$K$, if $X$ is understood) if $L$ is a separable
algebra over~$K$, \ie a direct product of separable
extensions $K_i$, and if the normalization $Y'$ of~$Y$
in $L$ (the disjoint sum of the normalizations of~$Y$
in the $K_i$) is unramified ($=$ \'etale by~\ref{I.9.11})
over~$Y$. Thus:

\begin{corollary}
\label{I.10.3}
For every $X$ finite over~$Y$ whose every irreducible
component dominates $Y$, let $R(X)$ be the ring of
rational functions on~$X$ (the product of the local
% original p. 22
rings of the generic points of the irreducible
components of~$X$), so that $X\mto R(X)$ is a functor
with values in the algebras of finite rank over
$K=R(Y)$. This functor establishes an equivalence of
the category of connected \'etale coverings of~$Y$
with the category of extensions $L$ of $K$ unramified
over~$Y$.
\end{corollary}

The inverse functor is the normalization functor.

Suppose that $Y$ is affine, hence defined by a normal
ring $A$ with field of fractions~$K$. Let $L$ be a
finite extension of~$K$ which is a direct product of
fields; then, by definition, the normalization $Y'$
of $Y$ in $L$ is isomorphic to $\Spec(A')$, where $A'$
is the normalization of~$A$ in~$L$. To say that $L$
is unramified over~$Y$ means that $A'$ is unramified
(or equivalently, \'etale) over~$A$. If $A$ is local,
this amounts to saying that the local rings
$A^\prime_{\goth{n}}$ (where $\goth{n}$ runs through
the finite set of maximal ideals of~$A'$, \ie its
prime ideals inducing the maximal ideal $\goth{m}$
of~$A$) are unramified ($=$ \'etale) over the local
ring~$A$. Finally, note also that the discriminant
criterion (\ref{I.4.10}) can also be applied in this
situation (more generally, a variant of the said
criterion should be stated as follows, without a
preliminary flatness condition when $X$ dominates
$Y$, with $Y$ nevertheless assumed locally integral:
$A\to B$ and $B\to B\otimes_A K=L$ are injective---then
$\trace_{L/K}$ is defined---and $\trace_{L/K}(xy)$
induces a \emph{fundamental bilinear form}
$B\times B\to A$, \ie there exist $x_i\in B$
($1\le i\le n$, $n=$ rank of~$L$ over~$K$) such that
$\trace(x_ix_j)\in A$ for every $i$,~$j$ and
$\det(\trace(x_ix_j))$ is invertible in~$A$).

The sorites (4.6) immediately implies the sorites
of unramifiedness in the classical setting:

\begin{proposition}
\label{I.10.4}
Let $Y$ be a normal integral prescheme, with field~$K$.
\textup{(i)} $K$ is unramified over~$Y$.
\textup{(ii)} If $L$ is an extension of~$K$ unramified
over~$Y$, if $Y'$ is a normal prescheme with field $L$
dominating $Y$ (for example, the normalization of~$Y$
in~$L$), and $M$ an extension of~$L$ unramified
over~$Y'$, then $M/K$ is unramified over~$X$
(\emph{transitivity} of unramifiedness).
\textup{(iii)} Let $Y'$ be a normal integral prescheme
dominating $Y$, with field $K'/K$; if $L$ is an
extension of~$K$ unramified over~$Y$, then
$L\otimes_K K'$ is an extension of~$K'$ unramified
over~$Y'$ (the \emph{translation} property).
\end{proposition}

Moreover:
% original p. 23

\begin{corollary}
\label{I.10.5}
Under the conditions of \textup{(iii)}, if
$Y=\Spec(A)$, $Y'=\Spec(A')$, then the normalization
$\bar{A'}$ of~$Y'$ in $L'=L\otimes_K K'$ is identified
with $\bar{A}\otimes_A A'$, where $\bar{A}$ is the
normalization of $A$ in $L$.
\end{corollary}

Usually, people (who are reluctant to consider rings
which are not integral domains, even if they are
direct products of fields) state the translation
property in the following (weaker) form:

\begin{corollary}
\label{I.10.6}
Under the conditions of \textup{(iii)}, let $L_1$ be
a \emph{compositum} of~$L/K$ (unramified over~$Y$)
and~$K'/K$. Then $L_1/K'$ is unramified over~$Y'$.
When $Y=\Spec(A)$, $Y'=\Spec(A')$, one moreover has
\[
\overline{A'}=A[\overline{A},A']
\]
\ie the ring $\overline{A'}$, the normalization
of~$A'$ in $L_1$, is the $A$-algebra generated by
$A'$ and the normalization $\overline{A}$ of~$A$
in $L$.
\end{corollary}

This last fact is false without the unramifiedness
assumption, even in the case of composita of number
fields\ldots

To conclude this no., we shall give the interpretation
of the notion of \'etale covering corresponding to
its intuitive picture: there should be the ``maximum
number'' of points above the point $y\in Y$ under
consideration, and in particular there should not be
``several coincident points'' above~$y$. To prove
results in this direction with all the desired
generality, we shall admit here Proposition~\ref{I.10.7}
below (whose proof will appear in the multiplodoque,
Chap\ptbl IV, par\ptbl 15, and uses Chevalley's
technique of constructible sets, and a little descent
theory\ldots).

A morphism of finite type $f\colon X\to Y$ is said
to be \emph{universally open} if for every extension
of the base $Y'\to Y$ (with $Y'$ locally noetherian)
the morphism $f'\colon X'=X\times_Y Y'\to Y'$ is
open, \ie sends open sets to open sets. One can,
moreover, restrict oneself to the case where $Y'$
is of finite type over~$Y$ (and even where $Y'$ is
of the form $Y[t_1,\dots,t_r]$, with the $t_i$
indeterminates). A universally open morphism is
a fortiori open (the converse being false);
on the other hand, if $f$ is open, with $X$ and $Y$
irreducible, then all the components of all the
fibers of~$f$ have the same dimension (namely, the
dimension of the generic fiber
% original p. 24
$f^{-1}(z)$, with $z$ the generic point of~$Y$).
Finally, if $Y$ is normal, this last condition
already implies that $f$ is \emph{universally} open
(Chevalley's theorem). It follows, for example,
that if $f\colon X\to Y$ is a \emph{quasi-finite}
morphism, with $Y$ normal and irreducible, then $f$
is universally open (or equivalently, open) if and
only if every irreducible component of~$X$ dominates~$Y$.
Recall also that a flat morphism (of finite type),
being open, is also universally open. With these
preliminaries in place, let us ``recall'' the following

\begin{proposition}
\label{I.10.7}
Let $f\colon X\to Y$ be a quasi-finite separated
universally open morphism. For every $y\in Y$, let
$n(y)$ be the ``geometric number of points of the
fiber $f^{-1}(y)$'', equal to the sum of the separable
degrees of the residue field extensions
$\kres(x)/\kres(y)$, for the points $x\in f^{-1}(y)$.
Then the function $y\mto n(y)$ on~$Y$ is upper
semicontinuous. For it to be constant in a neighborhood
of the point $y$ (\ie for one to have $n(y)=n(z_i)$,
where the $z_i$ are the generic points of the
irreducible components of~$Y$ containing $y$), it is
necessary that there exist a neighborhood $U$ of~$y$
such that $X|U$ is \emph{finite} over~$U$.\footnote{\Cf
EGA IV 15.5.1.}
\end{proposition}

\begin{corollary}
\label{I.10.8}
If $y\mto n(y)$ is constant and $Y$ is geometrically
unibranch\footnote{For the definition, \cf below,
\No\ref{I.11}, p\ptbl\pageref{I.11}.}, the irreducible
components of~$X$ are disjoint.
\end{corollary}

\begin{proposition}
\label{I.10.9}
Let $f\colon X\to Y$ be an \emph{\'etale} separated
morphism. With the notation of~\ref{I.10.7}, the
function $y\mto n(y)$ is upper semicontinuous. For
it to be constant in a neighborhood of the point $y$
(\ie for one to have $n(y)=n(z_i)$, where the $z_i$
are the generic points of the irreducible components
of~$Y$ containing $y$), it is necessary and sufficient
that there exist an open neighborhood $U$ of~$y$
such that $X|U$ is finite over U, \ie is an
\emph{\'etale covering} of~$U$.
\end{proposition}

\begin{corollary}
\label{I.10.10}
For a separated \'etale morphism $f\colon X\to Y$,
with $Y$ connected, to be finite (\ie to make $X$
an \emph{\'etale covering} of~$Y$), it is necessary
and sufficient that all the fibers of~$f$ have the
same geometric number of points.
\end{corollary}

In~\ref{I.10.7} and its corollary, there was no
normality assumption on~$Y$. If one makes such an
assumption, one obtains the stronger statement
(most often taken as the definition of a covering
being unramified):

\begin{theorem}
\label{I.10.11}
Let
% original p. 25
$f\colon X\to Y$ be a \emph{quasi-finite} separated
morphism. Suppose that~$Y$ is \emph{irreducible},
that every component of~$X$ dominates $Y$, and that
$X$ is reduced (\ie $\cal{O}_X$ has no nilpotent
elements). Let $n$ be the degree of~$X$ over~$Y$
(the sum of the degrees, over the field $K$ of~$Y$,
of the fields $K_i$ of the irreducible components
$X_i$ of~$X$). Let $y$ be a normal point of~$Y$.
Then the geometric number $n(y)$ of points of~$X$
above~$y$ is $\le n$, with equality if and only if
there exists an open neighborhood $U$ of~$y$ such
that $X|U$ is an \emph{\'etale covering} of~$U$.
\end{theorem}

The ``only if'' being trivial, let us prove the ``if''.
Let $z$ be the generic point of~$Y$; one has $n(z)=$
(sum of the separable degrees of the $K_i/K$) $\le n$,
and by~\ref{I.10.7}, $n(y)\le n(z)\le n$, with equality
implying that $X|U$ is \emph{finite} over~$U$ for a
suitable neighborhood~$U$ of~$y$. One may therefore
suppose that $X$ is finite over~$Y$ and that the
function $n(y')$ on~$Y$ is constant. Finally,
by~\ref{I.10.8}, $X$ is then the disjoint union of
its irreducible components, and to prove that it
is unramified at $y$, one is reduced to the case
where $X$ is irreducible, hence integral. Finally,
one may suppose that $Y=\Spec(\cal{O}_y)$. The theorem
then reduces to the following classical statement:

\begin{corollary}
\label{I.10.12}
Let $A$ be a normal local ring (noetherian, as always)
with field $K$, $L$ a finite extension of~$K$ of
degree $n$ and separable degree $n_s$, $B$ a subring
of~$L$ finite over~$A$, with field of fractions $L$,
$\goth{m}$ the maximal ideal of~$A$, and $n'$ the
separable degree of~$B/\goth{m}B$ over~$A/\goth{m}A=k$
($=$ the sum of the separable degrees of the residue
field extensions of this ring). One has $n'\le n_s$
and a fortiori $n'\le n$. This last inequality is
an equality if and only if $B$ is unramified
($=$ \'etale) over~$A$.
\end{corollary}

It remains only to show that $n'=n$ implies that $B$
is \'etale over~$A$. Let us recall the proof when $k$
is infinite: one need only show that $R=B/\goth{m}B$
is separable over~$k$; if this were not so, it would
follow (by a known lemma) that there exists an
element $a$ of~$R$ whose minimal polynomial over~$k$
has degree $>n'$. This element comes from an element
$x$ of~$B$, whose minimal polynomial over~$K$
(as an element of $L$) has degree $\le n$; on the
other hand, the latter has its coefficients in~$A$
since $A$ is normal, and hence gives, by reduction
mod $\goth{m}$, a monic polynomial $F\in k[t]$ of
degree $\le n=n'$ such that $F(a)=0$, a contradiction.

In
% original p. 26
the general case ($k$ possibly finite), returning
to geometric language, one considers $Y'=\Spec(A[t])$,
which is faithfully flat over~$Y$, and the generic
point $y'$ of the fiber $\Spec(k[t])$ of~$Y'$ over~$y$.
Then $X$ is net over~$Y$ at $y$ if and only if
$X'=X\times_Y Y'=\Spec(B[t])$ is net at $y'$ over~$Y'$,
as one sees at once. Moreover, by the choice of~$y'$,
its residue field is $k(t)$, hence infinite. Since
$y'$ is a normal point of~$Y'$, one is reduced to
the preceding case.

\section{Some complements}
\label{I.11}

We have already said that a connected \'etale covering
of an integral scheme is not necessarily integral.
Here are two examples of this fact.

a) Let $C$ be an algebraic curve with an ordinary double
point~$x$, $C^\prime$ its normalization, and $a$ and~$b$
the two points of~$C^\prime$ above~$x$. Let $C_i^\prime$
($i=1$, $2$) be two copies of~$C^\prime$, and $a_i$
and~$b_i$ the point of $C_i^\prime$ corresponding to $a$
\resp $b$. In the sum curve $C_1^\prime\amalg C_2^\prime$,
identify $a_1$ with~$b_2$, and $a_2$ with~$b_1$ (we leave
it to the reader to make the identification process
precise; it will be explained in Chap\ptbl VI of the
multiplodoque but, in the case of curves over an
algebraically closed field, is treated in Serre's book
on algebraic curves). One obtains a \emph{connected}
and \emph{reducible} curve $C^{\prime\prime}$ which is
an \'etale covering of degree~$2$ of~$C$. The reader
will verify that, in general, the connected ``Galois''
\'etale coverings $C^{\prime\prime}$ of~$C$ whose inverse
image $C^{\prime\prime}\times_C C^\prime$ is a
\emph{trivial} covering of~$C^\prime$ (\ie isomorphic
to the sum of a certain number of copies of~$C^\prime$)
are ``cyclic'' of degree~$n$, and that for every
integer~$n>0$ one can construct a connected cyclic
\'etale covering of degree~$n$. In the language of the
fundamental group, which will be developed later,
this means that the quotient of~$\pi_1(C)$ by the
closed normal subgroup generated by the image
of~$\pi_1(C^\prime)\to\pi_1(C)$ (the homomorphism
induced by the projection) is isomorphic to the
profinite completion of~$\ZZ$. More precisely, one
should be able to show that the fundamental group
of~$C$ is isomorphic to the (topological) free product
of the fundamental group of~$C^\prime$ with the
profinite completion of~$\ZZ$. Note that questions
of this kind gave rise to ``descent theory'' for schemes.

b) Let $A$ be a complete local integral domain. One
knows that its normalization $A^\prime$ is finite
% original p. 27
over~$A$ (Nagata), hence is a complete semilocal ring,
and therefore local since it is an integral domain.
Suppose that the residue field extension $L/k$ which
it defines is not purely inseparable (otherwise,
one will say that $A$ is geometrically unibranch;
\cf below). This is the case, for example, for the
ring $\RR[[s,t]]/(s^2+t^2)\RR[[s,t]]$, where $\RR$
is the field of real numbers. Let $k^\prime$ be a
finite Galois extension of~$k$ such that
$L\otimes_k k^\prime$ decomposes, and let~$B$ be
a finite \'etale algebra over~$A$ corresponding
to the residue field extension~$k^\prime$ (recall
that $B$ is essentially unique). Then
$B^\prime=A^\prime\otimes_A B$ over~$B$ has residue
algebra $L\otimes_k k^\prime$, which is not local;
hence $B^\prime$ is not a local ring, and therefore
(being complete) has zero divisors. Since it is
contained in the total ring of fractions of~$B$
(being free over~$A^\prime$, hence torsion-free
over~$A^\prime$, hence torsion-free over~$A$, and
therefore contained in
$B^\prime\otimes_A K=B^\prime_{(K)}
=A^\prime_{(K)}\otimes_K B_{(K)}=B_{(K)}$
since $A^\prime_{(K)}=K$), it follows that $B$ is
not an integral domain. In the case of the ring
$\RR[s,t]/(s^2+t^2)\RR[s,t]$, taking
$k^\prime/k=\CC/\RR$, one obtains for~$B$ the local
ring of two intersecting lines in the plane at
their point of intersection.

Note also that if there exists a connected \'etale
covering $X$ of an integral $Y$ which is not
irreducible, then every irreducible component of~$X$
gives an example of an unramified covering $X^\prime$
of~$Y$, dominating~$Y$, which is not \'etale over~$Y$.
In example~a), one thus finds that $C^\prime$ is
unramified over~$C$, without being \'etale at the
two points $a$ and~$b$ (as one also sees directly
by inspecting the completions of the local rings
of~$x$ and~$a$: from the ``formal'' point of view,
$C^\prime$ at the point~$a$ is identified with
a closed subscheme of~$C$ at the point~$x$, namely
one of the two ``branches'' of~$C$ through~$x$).

In a) and~b), one sees that the failure of the
conclusions of~\ref{I.9.5}\ptbl (i) and~(ii) is
directly linked to the fact that a point of~$Y$
``splits'' into \emph{distinct} points of the
normalization (in~b), the fact that the residue
field extension is not purely inseparable must
be interpreted geometrically in this way).
Precisely, we shall say that a local integral
domain~$A$ is \emph{geometrically unibranch} if
its normalization has only one maximal ideal,
with the corresponding residue field extension
purely inseparable; a point~$y$ of an integral
prescheme is said to be geometrically unibranch
if its local ring is so. Examples: a normal point,
an ordinary cusp of a curve, etc. It seems that
if $Y$ has a point which
% original p. 28
is not unibranch, there always exists a connected
\'etale covering of~$Y$ which is not irreducible;
at least this is what we have shown in case~b),
when $Y$ is the spectrum of a complete local ring.
One can show, on the other hand, that \emph{if
all points of~$Y$ are geometrically unibranch,
then every connected unramified $Y$-prescheme
dominating~$Y$ is \'etale} and irreducible.
The proof follows that of~\ref{I.9.5}, using the
following generalization of Theorem~\ref{I.8.3},
which will be proved later by means of descent
techniques\footnote{\Cf IX~\SourceRef{IX.4.10}{4.10}.
For a more direct proof, \cf EGA~IV 18.10.3,
using a variant of~\ref{I.9.5} for geometrically
unibranch local rings.}:

\emph{Let $Y^\prime\to Y$ be a finite, radicial,
surjective morphism (\ie what one might call
a ``universal homeomorphism''). Consider the
functor $X\mto X\times_Y Y^\prime=X^\prime$ from
$Y$-preschemes to $Y^\prime$-preschemes. This
functor induces an equivalence of the category
of \'etale $Y$-schemes with the category of
\'etale $Y^\prime$-schemes.} One can apply this
result, for example, when $Y^\prime$ is the
normalization of~$Y$, with $Y$ assumed unibranch
(and $Y^\prime$ finite over~$Y$, which holds in
all cases encountered in practice), or to a
$Y^{\prime\prime}$ ``sandwiched'' between~$Y$
and its normalization (which no longer needs
to be finite over~$Y$).

\end{document}
